\chapter{内部の仕組み}

本章では実装を開いて見ます。wick のゲームを書くために必要な知識ではありませんが、言語の設計は、その仕組みを通じて最もよく理解できます。wick の仕組みは、一章で率直に説明できるほど小さいものです。字句解析器、コンパイラ、仮想マシン、ガベージコレクタ、組み込み関数を含む言語全体が、次の三つのファイルにある、依存ライブラリのない約二千行の \mbox{C\texttt{++}} です。

\begin{center}
\begin{tabular}{@{}p{0.42\linewidth}p{0.50\linewidth}@{}}
\toprule
\texttt{wick/wick.hpp} & 組み込み用 API（約120行） \\
\texttt{wick/wick\_front.cpp} & 字句解析器と一回走査の型付きコンパイラ $\rightarrow$ バイトコード \\
\texttt{wick/wick\_vm.cpp} & スタック VM、マーク・スイープ GC、組み込み関数、シグネチャ解析器 \\
\bottomrule
\end{tabular}
\end{center}

\section{一回走査の型付きコンパイル}

特徴を決めている選択は、\textbf{抽象構文木がない}ことです。コンパイラは \emph{Crafting Interpreters} の \emph{clox} のような一回走査方式です。Pratt 法の式解析器が、進むにつれてバイトコードを生成し、式ごとに静的な `Type` を持ちます。型検査器とコード生成器が木を別々に二回たどるのではなく、同じ走査です。解析器が式を処理し終えると、その値を計算するバイトコードの生成と値の型の決定が、一度の走査でともに終わっています。

言語から見える二つの規則が、この設計から直接生まれます。

\paragraph{使用前の宣言。} コンパイラは名前や呼び出しに到達した時点で解決し、後から前方参照を補う走査はありません。そのため、関数は呼び出すより前に定義し、変数は使うより前に宣言しなければなりません。設計者が制約そのものを目的として選んだのではなく、一回走査のコンパイルの代償です。その代償は小さいものです。

\paragraph{生成時に型が分かること。} 式のコードを生成するとき、その型が分かっているので、コンパイラはバイトコードを\emph{特化}できます。性能の説明の中心なので、具体的に見ていきましょう。

\section{特化した命令}

動的な型の VM では、`a + b` は汎用の「加算」命令一つになります。実行時にオペランドのタグを見て、数値の加算、文字列の連結、エラーのどれを行うか決めます。そのタグ検査は、命令を実行するたびに行います。wick はすでに型を知っているため、コンパイル時に\emph{一度だけ}そのコストを払います。

\begin{itemize}
\item 二つの `num` への `+` は `OP_ADD`、二つの `str` への `+` は `OP_CONCAT` を生成します。選択はコード生成時に済んでいるので、実行時にタグを調べて選ぶことはありません。
\item 大小比較は、オペランドが `num` と分かっているため、数値比較命令 `OP_LT`、`OP_LE`、`OP_GT`、`OP_GE` にコンパイルします。
\item ローカル変数は\emph{フレームスロット}、つまり現在の呼び出しのスタックフレーム内への小さな整数オフセットに解決します。`OP_LOADL` と `OP_STOREL` はスロット番号を受け取り、実行時の名前検索はありません。グローバル変数も `OP_LOADG` と `OP_STOREG` でスロット添字に解決します。
\item エンジンのネイティブ関数はコンパイル時に整数 ID に解決します。VM のネイティブ呼び出し `OP_NCALL` は、名前でハッシュテーブルを探すのではなく、ネイティブ関数表への配列添字です。
\end{itemize}

\noindent
バイトコードには約四十の命令があります。算術と比較に加え、制御フローのジャンプ `OP_JMP`、`bool` を取り出す `OP_JF`、ループで後方へ飛ぶ `OP_JB`、オプショナル値を実装する二つのジャンプ `OP_JNN`（nil でなければジャンプ）と `OP_ISNIL` があります。呼び出し命令は wick 関数用の `OP_CALL` とネイティブ関数用の `OP_NCALL`、コンテナ命令は `OP_LIST`、`OP_MAP`、添字の取得・設定の組 `OP_IGET`・`OP_ISET` と `OP_MGET`・`OP_MSET` です。特化した組み込み関数には `OP_LEN`、`OP_TOSTR`、`OP_TONUM`、`OP_LPUSH`、`OP_LPOP` があります。いずれも、周囲のコードについて型検査器が証明した事実に対応します。

\section{スタック VM}

実行時の仕組みは、伝統的なスタックマシンです。定数プール、値のスタック、呼び出しフレームのスタックがあります。各フレームは、値スタック内でスロットが始まる基点と、命令ポインタを記録します。フレームを積むときには引数はすでにスタック上にあり、呼び出された関数のローカル変数はそのすぐ上のスロットに置かれます。戻るとスタックをフレームの基点まで縮めるので、トップレベルの文の間では空になり、特に大切なことに、フレームの間でも空になります。

実行時エラーには、コンパイル中に作るバイトコードオフセットと行番号の対応表を使います。範囲外の添字や失敗した `check` などのエラーを VM が出すとき、ソースの行を報告し、エンジンのエラー画面でその行を示せます。wick の実行時エラーが生のバイトコードオフセットではなく `file:line: message` になる理由です。

\textbf{JIT はありません}。これは未完成なのではなく、意図的です。400×240のゲームロジックにはインタープリタで十分な速さがあります。フレームごとの処理量は少なく、特化した命令が単純なインタープリタの余分なコストを取り除きます。さらに、エンジンは iOS をはじめ、実行時のコード生成と実行を禁止するプラットフォームを対象にしています。JIT があると、wick はそこで許可されません。特化したバイトコードのインタープリタは、仕組み上どこでも使え、ウォームアップや最適化解除による急な性能低下がなく、性能を予測できます。

\section{フレーム境界でのガベージコレクション}

wick は文字列、リスト、マップという三種類のヒープオブジェクトを、マーク・スイープ方式のコレクタで管理します。汎用のものと違い、ゲーム向けにしているのは、\emph{いつ}動くかです。自発的には決して動かず、ホストが `wick::collect()` を呼ぶときだけ動きます。lantern は画像を表示した後、毎フレーム正確に一度呼びます。

そのため、回収が描画の途中に割り込むことはありません。フレームが画面に出てから次が始まるまでという、停止が見えないタイミングだけで回収します。回収の最悪時のコストは、一フレームが作ったゴミの量で抑えられます。状態を再利用するゲームループでは、その量は小さいものです。Lua のゲームでフレームを落とすような、世代の状態に応じて予測できないタイミングで起きる停止はありません。

コレクタのルートは単純で、それは言語が小さい\emph{からこそ}です。マーク処理は定数プール、グローバル変数、VM の値スタックから始まり、リストとマップ内の参照を追います。v0.3 にはクロージャもアップバリューもないため、捕捉した環境をたどる必要がなく、オブジェクトグラフは浅い構造です。また、値スタックはフレーム間で空になるため、回収時に生きているルートはグローバル変数と定数、つまり長く使う状態だけで、マーク処理の仕事は少なくなります。

\section{wick の組み込み}

wick は組み込むために作られています。公開 API は `wick.hpp` だけで、組み込み側が内部ヘッダに触れる必要はありません。ライフサイクルは、作成、環境の登録、プログラムの読み込み、毎フレームの関数呼び出し、安全な時点での回収、ホットリロードのためのリセットです。

\begin{lstlisting}[language=C++,basicstyle=\ttfamily\small,keywordstyle=\color{wkKeyword}]
#include "wick.hpp"

wick::VM* vm = wick::create();
std::string err;

// typed natives: a signature DSL the compiler enforces
wick::addNative(vm, "game", "spawn(num, num, str): num",
    [](wick::VM& vm, const wick::Value* a, int) {
        int id = mySpawn(a[0].d, a[1].d, wick::getStr(a[2]));
        return wick::Value::num(id);
    }, err);
wick::addConst(vm, "game", "MAX", 64);

wick::load(vm, source, "main.wick", err);   // compile + run top level
wick::call(vm, "update", dt, true, err);    // each frame
wick::collect(vm);                          // at YOUR safe point
wick::reset(vm);                            // hot-reload support
\end{lstlisting}

\paragraph{シグネチャ DSL。} ネイティブ関数にはシグネチャ文字列を登録します。コンパイラは一度解析し、すべての呼び出し箇所で強制します。文法は次のとおりです。

\begin{quote}
\ttfamily name(type[, type][, type = default]) [: type]
\end{quote}

\noindent
型は `num`、`bool`、`str`、`list` で、末尾の `?` はオプショナル値を表します。この一つの文字列によって、`lt.load_save(str): str?` は保存の読み込みがオプショナル値を返すとコンパイラに教えます。エンジン呼び出しに対する第4章のオプショナル値の規律全体を、シグネチャ内の二文字が支えています。既定値は数値リテラルでなければならず、末尾の引数を呼び出し箇所で省略可能にします。コンパイラが補うため、ネイティブ関数は常にすべての引数を受け取り、`argc` を調べる必要がありません。名前空間 `"lt"` や `"game"` の下に登録すると、関数名が修飾されます。名前空間に null を渡すとグローバルスコープに置きます。`floor` や `sqrt` のような組み込み関数の登録方法です。

\paragraph{ネイティブ関数からのエラー報告。} ネイティブ関数は `wick::setError(vm, msg)` を呼んで失敗を知らせます。VM は呼び出し箇所を示す `file:line` の実行時エラーに変換するので、他の実行時エラーと同じ報告になり、同じエラー画面に出ます。値の補助関数 `makeStr`、`getStr`、`listLen`、`listGet` は、ネイティブ関数が受け取り返すヒープ上の値の読み取りと構築を担います。

\paragraph{ホットリロード。} `wick::reset(vm)` は、グローバル変数、関数、コンパイル済みプロトタイプというすべてのプログラム状態を消去し、登録済みの環境であるネイティブ関数と定数は残します。プログラムを解体し、バインディングを維持し、新しいソースを読み込む。ホットリロードに必要なことそのものです。`reset` は新しいプログラムを読み込む前に古いリソースを解放するので、再読み込みはリークしません。チュートリアルで約束した性質です。

どれも大きな仕組みではありません。組み込み API は約120行で、実装全体は約二千行です。小ささは本書を通じて繰り返す主題ですが、ここで最もはっきり見えます。静的な型を持ち、オプショナル値を検査し、決定的に動くスクリプト言語に、ガベージコレクタと型付きの外部関数インターフェースが備わっています。それでも午後だけで読み切れ、全体を自分で把握できます。
